
\item
\emph{Les parties \(\Lambda\) de \(\Pic_{X/S}\) correspondent biunivoquement aux familles \(\mathscr L\) de classes de faisceaux inversibles sur les fibres de \(X/S\).}
\item \emph{Une famille \(\Lambda\) de \(\Pic_{X/S}\) est quasi-compacte si et seulement si la famille \(\mathscr L\) correspondante est limit\'ee.}
\end{enumerate}

En effet, (i) r\'esulte de (3.1)\,; (ii), du fait suivant~: si \(s\) est un point de \(S\) et \(K\) est une extension alg\'ebriquement close de \(k(s)\), alors les \(K\)-points de \(\Pic_{X/S}\) correspondent biunivoquement aux classes d'isomorphisme de faisceaux inversibles sur \(X_K\).

Quant \`a (iii), supposons que la partie \(\Lambda\) de \(\Pic_{X/S}\) soit quasi-compacte. Alors \(\Lambda\) est contenue dans un ouvert \(U\) de type fini, et l'inclusion \(U\to\Pic_{X/S}\) correspond \`a un faisceau inversible \(L\) sur un \(X_T\), o\`u \(T\to U\) est un morphisme surjectif, (plat) de type fini. Enfin \(T\to S\) est de type fini, donc \(L\) limite \(\mathscr L\).

R\'eciproquement supposons que \(\mathscr L\) soit limit\'ee par un faisceau coh\'erent \(L\) sur un \(X_T\), avec \(T\) de type fini sur \(S\). On peut supposer \(L\) inversible (1.13 (i)) afin que \(L\) d\'efinisse un morphisme \(f:T\to\Pic_{X/S}\). L'image de \(f\) contient \(\Lambda\) et elle est quasi-compacte, donc noeth\'erienne par (i). Par suite \(\Lambda\) est quasi-compact et on a (iii).

Gr\^ace \`a 3.2, la d\'efinition habituelle des morphismes de type fini entre deux sch\'emas de Picard se g\'en\'eralise en la suivante~:

\textbf{D\'efinition 3.3.} \emph{Soient \(X\) et \(Y\) propres sur \(S\) noeth\'erien. Un morphisme entre les foncteurs de Picard}
\[
\begin{tikzcd}
\underline{\Pic}_{Y/S}\arrow[r]&\underline{\Pic}_{X/S}
\end{tikzcd}
\]
